\documentclass[12pt,a4paper]{article}

% ------------------------------------------------------------
% PACOTES
% ------------------------------------------------------------

\usepackage[utf8]{inputenc}
\usepackage[T1]{fontenc}
\usepackage[brazil]{babel}

\usepackage{times}
\usepackage{amsmath}
\usepackage{amssymb}
\usepackage{bm}

\usepackage{graphicx}
\usepackage{float}
\usepackage{booktabs}
\usepackage{array}
\usepackage{multirow}

\usepackage{cite}
\usepackage{url}

\usepackage{geometry}
\geometry{
    a4paper,
    left=3cm,
    right=2cm,
    top=3cm,
    bottom=2cm
}

\usepackage{setspace}
\usepackage{indentfirst}

\usepackage{caption}
\usepackage{subcaption}

\usepackage{listings}

% ------------------------------------------------------------
% CONFIGURAÇÕES
% ------------------------------------------------------------

\setlength{\parindent}{1.25cm}
\setlength{\parskip}{0pt}

\onehalfspacing

% ------------------------------------------------------------
% INFORMAÇÕES DO TRABALHO
% ------------------------------------------------------------

\title{
    \textbf{Controle e rastreamento de sistemas farmacocinéticos}
}

\author{
    Melvin Gustavo Maradiaga Elvir
}

\date{
    Campinas\\
    2026
}

% ------------------------------------------------------------
% DOCUMENTO
% ------------------------------------------------------------

\begin{document}

% FOLHA DE ROSTO
% ============================================================

\begin{titlepage}

\begin{center}

{\large
UNIVERSIDADE ESTADUAL DE CAMPINAS
}

\vspace{0.3cm}

{\large
FACULDADE DE ENGENHARIA ELÉTRICA E DE COMPUTAÇÃO
}

\vfill

{\large
MELVIN GUSTAVO MARADIAGA ELVIR
}

\vfill

{\Large
\textbf{Controle e rastreamento de sistemas farmacocinéticos}
}

{\large
<Link do video>
}


\vspace{1cm}

\begin{flushright}
\begin{minipage}{0.55\textwidth}
Trabalho apresentado como requisito para
conclusão do curso de Engenharia Elétrica,
sob orientação da Prof(a). Dr(a). Cecília de
Freitas Morais.
\end{minipage}
\end{flushright}

\vfill

Campinas\\
2026

\end{center}

\end{titlepage}

% ============================================================
% DEDICATÓRIA
% ============================================================

\section*{Dedicatória}

% Opcional.

\newpage

% ============================================================
% AGRADECIMENTOS
% ============================================================

\section*{Agradecimentos}

% Escrever posteriormente.

\newpage

% ============================================================
% RESUMO
% ============================================================

\section*{Resumo}

% Escrever posteriormente.
%
% O resumo deve apresentar:
% - contexto/problema;
% - objetivo;
% - metodologia;
% - principais resultados;
% - conclusão.
%
% Evitar referências bibliográficas no resumo.

\vspace{0.5cm}

\noindent
\textbf{Palavras-chave:} controle robusto; desigualdades matriciais
lineares; farmacocinética; sistemas ciberfísicos médicos;
sistemas embarcados.

\newpage

% ============================================================
% ABSTRACT
% ============================================================

\section*{Abstract}

% Escrever posteriormente.

\vspace{0.5cm}

\noindent
\textbf{Keywords:} robust control; linear matrix inequalities;
pharmacokinetics; medical cyber-physical systems;
embedded systems.

\newpage

% ============================================================
% SUMÁRIO
% ============================================================

\tableofcontents

\newpage

% ============================================================
% 1. INTRODUÇÃO
% ============================================================

\section{Introdução}

Os sistemas ciberfísicos \cite{lee2017embedded} são uma integração da computação com processos físicos, trabalhando
na interseção entre componentes cibernéticos, como os microcontroladores, e elementos que
interagem com o mundo físico, como sensores e atuadores. Eles têm sido implementados no
desenvolvimento de diversos projetos, variando em complexidade, até chegar a sistemas
inovadores, como veículos autônomos, infraestruturas críticas inteligentes e, nos últimos anos,
sistemas médicos ciberfísicos (MCPS) \cite{dey2018medical}. Com o avanço da biomedicina, o interesse no
desenvolvimento de MCPS tem crescido, impulsionando a implementação de tecnologias de controle
para ``fechar a malha'' (realimentar o sistema com sinais fisiológicos) e criar dispositivos
auto-regulados, como o pâncreas artificial \cite{brown2019six}, e controladores para reanimação volêmica em casos
de hipovolemia \cite{kramer2008closed}. Dentro deste contexto, a administração automatizada de medicamentos, na qual
algoritmos de controle são empregados para ajustar continuamente a dosagem administrada ao
paciente, se torna uma das aplicações mais promissoras dos MCPS.

Nesse cenário, o campo da farmacologia clínica é de especial interesse devido ao potencial que os
MCPS possuem para aumentar a eficácia terapêutica e diminuir os riscos associados à
administração inadequada de medicamentos. A área de farmacocinética, que estuda a concentração
de medicamentos no organismo em função de tempo e dose \cite{bailey2005drug}, facilita o desenvolvimento de tais
dispositivos através dos modelos farmacocinéticos compartimentais. Baseando-se neles, diversas
estratégias de controle têm sido propostas \cite{gholami2013optimal,nogueira2016positive,padmanabhan2019optimal}
para administração automatizada de medicamentos.
Uma alternativa de projeto para garantir a automatização do rastreamento de referências em
administração farmacológica seria adotar técnicas de controle robusto baseadas em programação
convexa, mediante o uso de desigualdades matriciais lineares (LMIs), para o projeto de
controladores, como feito em \cite{marcomini2025projeto}. A vantagem de empregar LMIs é a possibilidade de tratar sistemas
afetados por incertezas paramétricas e variantes no tempo, ou seja, o rastreador projetado seria
válido mesmo que a modelagem fisiológica do indivíduo não fosse precisa ou sofresse pequenas
variações ao longo do tempo. Contudo, a validação das estratégias de controle tem sido realizada
predominantemente em ambientes de simulação. Para aproximar os resultados obtidos de
condições reais de operação, plataformas físicas de emulação têm sido desenvolvidas para a
validação dos controladores, tanto através de circuitos eletrônicos \cite{costa2024designing} quanto por meio de sistemas
embarcados \cite{desouza2025embedded}.


% ============================================================
% 2. OBJETIVOS
% ============================================================

\section{Objetivos}

\subsection{Objetivo geral}

% Formular aqui o objetivo geral do trabalho.


\subsection{Objetivos específicos}

% Listar os objetivos específicos.

\begin{itemize}

    \item Selecionar e implementar um modelo farmacocinético
    adequado para representar a dinâmica de absorção e
    distribuição do fármaco;

    \item Formular o problema de rastreamento de referência
    considerando incertezas paramétricas;

    \item Projetar um controlador/rastreador robusto utilizando
    condições baseadas em desigualdades matriciais lineares
    (LMIs);

    \item Validar o controlador por meio de simulações
    computacionais em Matlab/Simulink;

    \item Implementar o modelo da planta em uma plataforma
    embarcada;

    \item Implementar o controlador em um segundo
    microcontrolador;

    \item Desenvolver a comunicação entre os dispositivos;

    \item Realizar a emulação da malha fechada;

    \item Comparar os resultados obtidos experimentalmente
    com os resultados obtidos por simulação;

    \item Avaliar o comportamento do sistema diante de
    diferentes referências e condições de incerteza.

\end{itemize}


% ============================================================
% 3. FUNDAMENTAÇÃO TEÓRICA
% ============================================================

\section{Fundamentação Teórica}

% Este capítulo apresenta os conceitos necessários para que
% o leitor compreenda o desenvolvimento do trabalho.
%
% Evitar apresentar aqui resultados próprios.
% A ideia é construir a base teórica utilizada posteriormente.


\subsection{Sistemas ciberfísicos médicos}

% Definição de CPS e MCPS.
%
% Relacionar computação, sensores, atuadores e processos físicos.
%
% Apresentar exemplos de aplicações médicas de controle
% em malha fechada.


\subsection{Farmacocinética e modelos compartimentais}

% Explicar:
% - conceito de farmacocinética;
% - concentração do fármaco;
% - compartimentos;
% - entrada de medicamento;
% - dinâmica de absorção/distribuição;
% - modelo utilizado no trabalho.


\subsection{Modelagem em espaço de estados}

% Apresentar:
%
% \begin{equation}
%     \dot{x}(t) = Ax(t) + Bu(t)
% \end{equation}
%
% \begin{equation}
%     y(t) = Cx(t) + Du(t)
% \end{equation}
%
% Explicar o significado das matrizes e estados.


\subsection{Controle e rastreamento de referência}

% Introduzir:
% - controle em malha fechada;
% - estabilidade;
% - erro de rastreamento;
% - referência;
% - realimentação.


\subsection{Incertezas paramétricas}

% Explicar por que os parâmetros do modelo farmacocinético
% podem apresentar incertezas.
%
% Definir como essas incertezas serão representadas
% matematicamente no trabalho.


\subsection{Desigualdades matriciais lineares}

% Apresentar o conceito de LMI.
%
% Explicar sua utilização na formulação de problemas
% de controle convexo.
%
% Apresentar somente a teoria necessária para compreender
% o projeto desenvolvido posteriormente.


\subsection{Controle robusto}

% Explicar o conceito de robustez.
%
% Relacionar:
%
% incerteza do modelo
%       ↓
% conjunto de modelos possíveis
%       ↓
% controlador robusto
%       ↓
% garantia de estabilidade/desempenho.


\subsection{Implementação embarcada e Hardware-in-the-Loop}

% Apresentar o conceito de Hardware-in-the-Loop utilizado
% neste trabalho.
%
% Explicar a separação entre:
% - dispositivo responsável pela planta;
% - dispositivo responsável pelo controlador;
% - comunicação entre os dispositivos.


% ============================================================
% 4. METODOLOGIA
% ============================================================

\section{Metodologia}

% Este capítulo deve explicar COMO o trabalho foi realizado.


\subsection{Modelo farmacocinético utilizado}

% Apresentar o modelo escolhido.
%
% Equações diferenciais.
% Estados.
% Entradas.
% Saídas.
% Parâmetros.


\subsection{Formulação do problema de controle}

% Definir matematicamente:
% - planta;
% - referência;
% - erro;
% - variável de controle;
% - requisitos de desempenho.


\subsection{Modelagem das incertezas}

% Mostrar como os parâmetros incertos serão representados.


\subsection{Projeto do rastreador robusto}

% Apresentar a estratégia de projeto.
%
% Não apenas mostrar o resultado:
% explicar a sequência matemática que leva às LMIs.


\subsection{Formulação das condições LMIs}

% Apresentar as LMIs utilizadas no projeto.
%
% Mostrar as variáveis de decisão.
% Mostrar as restrições.
% Explicar como a solução é obtida.


\subsection{Simulação computacional}

% Descrever:
% - Matlab;
% - Simulink, se utilizado;
% - condições iniciais;
% - referências;
% - valores dos parâmetros;
% - cenários de incerteza;
% - métricas avaliadas.


\subsection{Arquitetura da plataforma embarcada}

% Inserir uma figura da arquitetura geral.
%
% Exemplo:
%
% Microcontrolador 1
%      ↓
% Emulação da planta
%      ↓
% Comunicação
%      ↓
% Microcontrolador 2
%      ↓
% Controlador
%      ↓
% Ação de controle
%      ↺


\subsection{Implementação nos microcontroladores}

% Descrever:
% - microcontroladores utilizados;
% - linguagem;
% - bibliotecas;
% - frequência de execução;
% - representação numérica;
% - comunicação.


\subsection{Procedimento de validação}

% Explicar como os experimentos foram realizados.
%
% Definir os cenários de teste e os critérios utilizados
% para comparação com a simulação.


% ============================================================
% 5. DESENVOLVIMENTO E IMPLEMENTAÇÃO
% ============================================================

\section{Desenvolvimento e Implementação}

% Aqui começa efetivamente o desenvolvimento do seu projeto.


\subsection{Implementação do modelo da planta}

% Mostrar como o modelo matemático foi transformado
% em algoritmo executável.


\subsection{Implementação do controlador}

% Mostrar como o controlador obtido pelas LMIs foi
% implementado na plataforma embarcada.


\subsection{Implementação da comunicação}

% Descrever o protocolo e a estrutura dos dados trocados.


\subsection{Integração da malha fechada}

% Mostrar o sistema completo funcionando.
%
% Planta → comunicação → controlador → ação de controle.


\subsection{Plataforma experimental}

% Fotos, diagrama elétrico, arquitetura e características
% do hardware.


\subsection{Execução dos ensaios}

% Descrever os testes efetivamente realizados.


% ============================================================
% 6. RESULTADOS E DISCUSSÃO
% ============================================================

\section{Resultados e Discussão}


\subsection{Resultados da síntese do controlador}

% Mostrar:
% - solução das LMIs;
% - ganhos obtidos;
% - parâmetros relevantes;
% - condições de projeto.


\subsection{Resultados de simulação}

% Gráficos:
% - referência × saída;
% - erro;
% - ação de controle;
% - estados;
% - diferentes condições de incerteza.


\subsection{Resultados da implementação embarcada}

% Apresentar os resultados obtidos no hardware.


\subsection{Comparação entre simulação e emulação}

% Comparar quantitativamente:
%
% Simulação × Hardware
%
% Erro de rastreamento
% Tempo de resposta
% Sobressinal, se aplicável
% Ação de controle
% Comportamento diante das incertezas


\subsection{Análise da robustez}

% Testar diferentes valores dos parâmetros incertos.
%
% Mostrar se o comportamento esperado do controlador
% permanece válido.


\subsection{Limitações e considerações práticas}

% Discutir:
% - limitações do modelo;
% - limitações do hardware;
% - limitações da comunicação;
% - limitações computacionais;
% - diferenças entre simulação e sistema embarcado.


% ============================================================
% 7. CONCLUSÃO
% ============================================================

\section{Conclusão}

% Retomar:
% - problema;
% - objetivo;
% - metodologia;
% - principais resultados;
% - contribuição do trabalho.
%
% Não introduzir resultados que não tenham sido apresentados
% anteriormente.


\subsection{Trabalhos futuros}

% Possíveis extensões:
% - comunicação remota;
% - implementação em hardware mais próximo de uma aplicação real;
% - novos modelos farmacocinéticos;
% - outras estratégias de controle;
% - testes com maior variedade de incertezas;
% - integração com sistemas de monitoramento.


% ============================================================
% REFERÊNCIAS
% ============================================================

\section{Referências}



% Inicialmente você pode inserir as referências manualmente.
%
% Posteriormente, recomendo migrar para BibTeX/BibLaTeX
% para facilitar a organização.


% ============================================================
% APÊNDICE
% ============================================================

\appendix

\section{Código do controlador}

% Inserir posteriormente, se necessário.


\section{Código da implementação embarcada}

% Inserir posteriormente, se necessário.


% ============================================================
% ANEXO
% ============================================================

\section{Materiais complementares}

% Documentos, tabelas, informações externas etc.


\end{document}